\section{完美域与非完美域}\label{sec:galois-2-6}
\subsection{完美域}\label{sec:galois-2-6-1}
\begin{definition*}
设域 $k$ 的特征为 $p$. 若 $p=0$，或其 Frobenius 自同态为双射，则称 $k$ 为\newterm{galois-perfect-field}{完美域}{perfect field}.
\end{definition*}
注意，特征为 $p>0$ 的域的 Frobenius 自同态 $F_k$ 总是单射（所有域态射都是单射），所以 $k$ 完美当且仅当 $F_k$ 为满射.
\begin{example*}
\begin{enumerate}[label=\roman*)]
\item 每个有限域都是完美域，因为有限集到自身的单射也是双射.
\item 每个代数闭域都是完美域，因为对每个 $x\in k$，方程 $X^p-x$ 都有解.
\item 域 $\mathbb F_p(T)$ 不是完美域，因为 $T$ 没有 $p$ 次根，故不属于 Frobenius 自同态的像.
\end{enumerate}
\end{example*}
\subsection{不可约多项式的可分性}\label{sec:galois-2-6-2}
设 $f\in k[X]$ 不可约. 由于 $(f)$ 是极大理想，可见 $(f,f')=k[X]$ 当且仅当 $f$ 不整除 $f'$. 由 $\deg(f')<\deg(f)$，这又等价于 $f'\ne0$. 因此
\begin{align*}
f\text{ 可分}\quad\Longleftrightarrow\quad f'\ne0.
\end{align*}
由此容易得到下面的定理.
\begin{theorem*}
设 $f\in k[X]$ 不可约.
\begin{enumerate}[label=\roman*)]
\item 若 $k$ 完美，则 $f$ 可分.
\item 若 $f$ 不可分，则 $k$ 不完美且特征为 $p>0$，并存在唯一的不可约且可分的 $g\in k[X]$ 及唯一的 $r\in\mathbb N$，使得 $f=g(X^{p^r})$.
\end{enumerate}
\end{theorem*}
\begin{proof}
i) 有 $\deg(f)>0$，所以若 $k$ 的特征为零，则 $f'\ne0$，由前述结论，$f$ 可分. 设 $k$ 为特征 $p>0$ 的完美域. 若 $f$ 不可分，则 $f'=0$，故 $f=g(X^p)$，其中 $g=\sum_na_nX^n\in k[X]$. 由于 $k$ 完美，可找到 $b_n$ 使 $a_n=b_n^p$. 置 $\widetilde g:=\sum_nb_nX^n$. 则 $f=\sum_nb_n^pX^{np}=\widetilde g^p$，与 $f$ 不可约矛盾.

ii)\footnote{译注：原文此证明段的分目误印为 iii)，对应定理的 ii).} 由于 $f$ 不可分，有 $f'=0$，故存在 $f_1$ 使得 $f=f_1(X^p)$. 多项式 $f_1$ 显然也不可约，因此可以递归重复这一过程：或者 $f_1$ 可分，或者存在 $f_2$ 使得 $f_1=f_2(X^p)$，如此继续. 由于次数下降，存在 $r$ 使得 $f_r$ 可分. 唯一性留给读者验证.
\end{proof}
\begin{corollary*}
完美域上的每个代数扩张都是可分扩张.
\end{corollary*}
\begin{remark*}
在非完美域上，存在不可分的不可约多项式. 事实上，取不属于 Frobenius 像的 $\alpha\in k$. 多项式 $f=X^p-\alpha$ 不可分. 下面证明它不可约. 设 $f=f_1f_2$ 是 $f$ 在 $k[X]$ 中的非平凡分解. 若 $\beta$ 是 $\alpha$ 在 $k$ 的代数闭包 $\overline k$ 中的一个 $p$ 次根，则在 $\overline k[X]$ 中有 $f=(X-\beta)^p$，从而 $f_i=(X-\beta)^{r_i}$，其中 $0<r_i<p$.\footnote{译注：此处原文未说明把因子 $f_i$ 取为首一；所写等式按首一因子理解.} 注意 $r_1$ 与 $r_2$ 互素，所以存在 $u,v$ 使得 $ur_1+vr_2=1$. 因而在 $\overline k(X)$ 中有等式 $f_1^uf_2^v=X-\beta$. 由于 $k(X)\cap\overline k[X]=k[X]$（交取在 $\overline k(X)$ 内），可知 $X-\beta\in k[X]$，从而 $\beta\in k$，与对 $\alpha$ 的假设矛盾.
\end{remark*}
这一注记说明，以下三个断言事实上等价：
\begin{itemize}
\item $k$ 完美.
\item 每个不可约多项式 $f\in k[X]$ 都可分.
\item $k$ 的每个代数扩张都可分.
\end{itemize}
这无疑解释了``完美''这一术语. 那么，对非完美域会发生什么？
\subsection{非完美域与完全不可分扩张}\label{sec:galois-2-6-3}
\begin{definition*}
设 $K\supset k$ 为代数扩张.
\begin{itemize}
\item 若元素 $\alpha\in K$ 的极小多项式具有 $X^{p^r}-x$ 的形式，其中 $x\in k$，则称 $\alpha$ 在 $k$ 上\newterm{galois-purely-inseparable-element}{完全不可分}{purely inseparable element}.
\item 若每个 $\alpha\in K\setminus k$ 在 $k$ 上完全不可分，则称 $K$ 在 $k$ 上完全不可分，称此扩张为\newterm{galois-purely-inseparable-extension}{完全不可分扩张}{purely inseparable extension}.
\end{itemize}
\end{definition*}
\begin{example*}
扩张 $K=\mathbb F_p(T)\supset\mathbb F_p(T^p)=k$ 完全不可分. 事实上，对每个 $\alpha\in K$ 都有 $\alpha^p\in k$，故 $f_\alpha$ 整除 $X^p-\alpha^p=(X-\alpha)^p$，当 $\alpha\notin k$ 时便与它相等.

更一般地，若 $k$ 为正特征域，且 $x\in k$ 不在 Frobenius 像中，则 $K:=k[X]/(X^p-x)$ 是完全不可分扩张.
\end{example*}
下面的命题从可分性的角度，描述非完美域上代数扩张的结构.
\begin{proposition*}
设 $K\supset k$ 为代数扩张. 则 $K$ 中在 $k$ 上可分的元素构成的集合 $K_{\mathrm{sep}}$ 是 $K$ 的子域，称为 $k$ 在 $K$ 中的\newterm{galois-relative-separable-closure}{可分闭包}{relative separable closure}，且扩张 $K\supset K_{\mathrm{sep}}$ 完全不可分.
\end{proposition*}
\begin{proof}
设 $\alpha,\beta\in K_{\mathrm{sep}}$. 扩张 $k\subset k[\alpha]$ 可分. $\beta$ 在 $k[\alpha]$ 上的极小多项式整除 $f_\beta$，故可分，从而扩张 $k[\alpha]\subset k[\alpha,\beta]$ 也可分. 由推论~\ref{sec:galois-2-5-4}，$k[\alpha,\beta]$ 在 $k$ 上可分，特别地，$\alpha-\beta$ 与 $\alpha\beta^{-1}$ 在 $k$ 上可分.\footnote{译注：原文未明写后式需 $\beta\ne0$.} 因此 $K_{\mathrm{sep}}$ 是域.

再取 $\alpha\in K\setminus K_{\mathrm{sep}}$，设 $g_\alpha$ 为其在 $K_{\mathrm{sep}}$ 上的极小多项式. 取 $r$ 使得 $g_\alpha=h_\alpha(X^{p^r})$，其中 $h_\alpha$ 可分. 则 $h_\alpha$ 的每个根都在 $K_{\mathrm{sep}}$ 上可分，从而也在 $k$ 上可分（仍由推论~\ref{sec:galois-2-5-4}），故属于 $K_{\mathrm{sep}}$.\footnote{译注：原文如此表述；由 $K_{\mathrm{sep}}$ 的定义，只能对属于 $K$ 的根直接得出此结论. 此处 $\alpha^{p^r}\in K$ 本身就是 $h_\alpha$ 的一个根，足以继续原论证.} 由于 $h_\alpha$ 不可约，有 $h_\alpha=X-\beta$，其中 $\beta\in K_{\mathrm{sep}}$. 因而 $g_\alpha=X^{p^r}-\beta$.
\end{proof}
\begin{definition*}
\begin{enumerate}[label=\roman*)]
\item 若 $k[X]$ 中每个不可约可分多项式都分裂，则称域 $k$ \newterm{galois-separably-closed}{可分闭}{separably closed}.
\item 域 $k$ 的\newterm{galois-separable-closure}{（绝对）可分闭包}{separable closure}是指 $k$ 的一个代数、可分且可分闭的扩张.
\end{enumerate}
\end{definition*}
\begin{proposition*}
每个域 $k$ 都有可分闭包，且在同构意义下唯一. 此外，每个可分扩张都可嵌入其中.
\end{proposition*}
\begin{proof}
设 $\overline k$ 为 $k$ 的一个代数闭包. 则 $\overline k_{\mathrm{sep}}$ 是 $k$ 的可分闭包（细节请自行验证）.
\end{proof}

